\providecommand{\todo}[1]{\textcolor{red}{\textbf{[TODO: #1]}}}
\clearpage
\phantomsection
\section*{Chương 7. Tự tương quan bằng MAC}
\addcontentsline{toc}{section}{Chương 7. Tự tương quan bằng MAC}

\begin{learningobj}
\begin{itemize}[leftmargin=1.5em, itemsep=2pt]
    \item Thấu hiểu bản chất toán học và cơ chế vật lý của hàm tự tương quan thời gian thực $R[\tau] = \sum_{n=0}^{N-1-\tau} x_w[n] \cdot x_w[n+\tau]$; dẫn xuất chính xác phép đếm số lượng phép tính nhân-cộng: $\sum_{\tau=26}^{214} (638 - \tau) = 97.902\text{ MAC}$ cho dải quét và $98.540\text{ MAC}$ khi tính kèm $R[0]$.
    \item Khảo sát mã nguồn Praat 7.0.03 (\path{fon/Sound_to_Pitch.cpp}, dòng 79--80, 144--164, 344--345, 401--403); giải mã định lý đệm số $0$ từ $638$ lên $1024\text{ điểm}$ FFT của Paul Boersma và chứng minh bằng giải tích rằng giới hạn không cuộn vòng (no circular wrap-around) đạt tới $\tau \le 386\text{ mẫu}$, bảo đảm tự tương quan vòng qua FFT đồng nhất tuyệt đối với tự tương quan tuyến tính qua MAC trực tiếp (sai số dấu phẩy động $\sim 10^{-16}$).
    \item Phân tích toàn diện sự đánh đổi giữa thuật toán FFT và kiến trúc MAC trực tiếp trên FPGA; chứng minh rằng với ngân sách $500.000\text{ chu kỳ/khung}$ ($10\text{ ms}$ tại $50\text{ MHz}$), một bộ MAC đơn lẻ trên 1 khối DSP chỉ chiếm $19{,}7\%$ thời gian thực ($\text{RTF} \approx 0{,}197$), loại bỏ hoàn toàn sự phức tạp của bộ đệm ping-pong, bộ quay pha twiddle ROM và mạch quản lý tăng trưởng bit.
    \item Khảo sát sâu sắc ba phương án chốt thanh ghi đường ống trong khối DSP Cyclone V: phân tích sự thật không thể bỏ thanh ghi tích lũy, so sánh giữa không chốt toán hạng (trễ 2 nhịp) và chốt toán hạng tại cửa khối DSP (trễ 3 nhịp, tận dụng thanh ghi cứng có sẵn 0 ALM); làm rõ sự khác biệt giữa độ trễ (latency) và thông lượng (throughput), chứng minh độ trễ 3 chu kỳ chỉ làm sai khác $0{,}001\%$ thời gian thực trên $98.543\text{ chu kỳ}$.
    \item Dẫn xuất độ rộng từ bit toán hạng $17\text{-bit}$ có dấu và chứng minh bằng toán học cận trên $43\text{-bit}$ của bộ cộng tích lũy, đính chính nhận định "48 bit tự nhiên" của Argon bằng việc chỉ rõ bộ tích lũy phần cứng Cyclone V rộng tới $64\text{ bit}$ (CV-52003 Chương 5).
    \item Thiết kế và hiện thực hóa hai đường dẫn phần cứng: Tuyến (A) cấu hình khối IP Intel \texttt{ALTERA\_MULT\_ADD} (UG-01069) cùng cờ nạp đè đồng bộ \texttt{accum\_sload} để chuyển tiếp lag zero-bubble; Tuyến (B) module HDL tự viết \texttt{ch07\_mac\_unit.v}, bộ đệm khung True Dual-Port RAM 2 khối M10K \texttt{ch07\_xw\_ram.v} và khối điều khiển phát địa chỉ liên tục \texttt{ch07\_mac\_engine.v}.
    \item Thực thi quy trình kiểm chứng chu-kỳ-đối-chu-kỳ giữa mô hình song sinh Python (\texttt{ch07\_mac\_twin.py}) và mô phỏng Verilog trên $7.460\text{ vector}$, tái hiện toàn bộ $3.040\text{ kết quả}$ đường chuẩn từ \texttt{golden\_v2/hdl\_check/} trên 16 khung tín hiệu bao gồm 2 khung stress biên cực hạn ($|x_w| = 59.379$).
\end{itemize}
\end{learningobj}

\subsection*{7.1. Bản chất của hàm tự tương quan và phép đếm số lượng phép tính MAC}
\addcontentsline{toc}{subsection}{7.1. Bản chất của hàm tự tương quan và phép đếm số lượng phép tính MAC}

\subsubsection*{Vấn đề thực tế: Sự lừa dối của miền thời gian và sự cần thiết của hàm tự tương quan}
Hãy xem xét bài toán cốt lõi mà kỹ sư xử lý tín hiệu phải đối mặt khi ước lượng cao độ giọng nói: ta cần tìm khoảng thời gian lặp lại cơ bản $T_0$ (tính theo số mẫu rời rạc) của dây thanh âm trong một đoạn tín hiệu hữu hạn dài $N = 638\text{ mẫu}$. Giả sử tín hiệu âm thanh thu được là nguyên âm /a/ có tần số cơ bản $F_0 = 150\text{ Hz}$ tại tần số lấy mẫu $f_s = 16.000\text{ Hz}$, tương ứng với chu kỳ danh định $\tau_0 = 16.000 / 150 \approx 106{,}67\text{ mẫu}$.

Nếu người thiết kế ngây thơ tìm kiếm chu kỳ bằng cách dò khoảng cách giữa các điểm giao cắt mức không (zero-crossing) hoặc khoảng cách giữa các đỉnh biên độ tức thời, thuật toán sẽ sụp đổ ngay lập tức. Dây thanh âm không phát ra sóng sin thuần khiết; nó phát ra một chuỗi xung tam giác giàu họa ba, kích thích các khoang cộng hưởng của thanh quản và vòm họng (formants). Các họa ba bậc hai ($300\text{ Hz}$), bậc ba ($450\text{ Hz}$) và bậc bốn ($600\text{ Hz}$) giao thoa sóng với nhau, tạo ra vô số đỉnh nhọn cục bộ và các điểm cắt qua số $0$ giả mạo bên trong một chu kỳ đơn lẻ. Tệ hơn nữa, sự dịch pha do bộ lọc âm học làm thay đổi hoàn toàn hình dáng trực quan của sóng âm trong miền thời gian: hai chu kỳ liên tiếp có cùng chu kỳ lặp lại nhưng đỉnh biên độ tức thời có thể lệch nhau hàng chục mẫu.

Làm thế nào để đo lường chu kỳ lặp lại của một dạng sóng phức tạp mà không bị đánh lừa bởi hình dáng cục bộ? Câu trả lời kinh điển của lý thuyết truyền thông là: \textbf{So sánh dạng sóng với chính bản sao bị dịch thời gian của nó}. Nếu ta dịch tín hiệu đi đúng một khoảng trễ $\tau$, tại $\tau = \tau_0$, toàn bộ cấu trúc họa ba và xung đỉnh của tín hiệu sẽ chồng khít lên bản sao gốc. Tích vô hướng giữa hai dạng sóng sẽ đạt cực đại toàn cục. Phép toán toán học thực hiện sự so khớp toàn diện này chính là \textbf{Hàm tự tương quan} (Autocorrelation Function).

\subsubsection*{Dẫn xuất từ nguyên lý gốc: Công thức tự tương quan hữu hạn}
Với một khung âm thanh đã nhân cửa sổ Hanning $x_w[n]$ có chiều dài hữu hạn $N = 638\text{ mẫu}$ ($n = 0, 1, \dots, N-1$), hàm tự tương quan ngắn hạn không chuẩn hóa tại độ trễ $\tau$ được định nghĩa theo nguyên lý tích chập năng lượng:
\begin{equation}
R[\tau] = \sum_{n=0}^{N-1-\tau} x_w[n] \cdot x_w[n+\tau]
\label{eq:ch07_autocorr_def}
\end{equation}

Mỗi giá trị $R[\tau]$ trong phương trình~(\ref{eq:ch07_autocorr_def}) là tổng của các tích giữa mẫu $x_w[n]$ tại thời điểm hiện tại và mẫu $x_w[n+\tau]$ tại thời điểm tương lai dịch đi $\tau$ bước. Hãy quan sát cận trên của tổng tích lũy:
\begin{itemize}[leftmargin=1.5em, itemsep=2pt]
    \item Khi $\tau = 0$: Mẫu dịch chuyển trùng khít với mẫu gốc. Ta có tổng $N = 638\text{ tích}$:
    \begin{equation}
    R[0] = \sum_{n=0}^{637} x_w^2[n]
    \end{equation}
    Đại lượng $R[0]$ chính là năng lượng tổng cộng của khung âm thanh sau cửa sổ hóa, đóng vai trò mốc chuẩn hóa biên độ và kiểm tra điều kiện im lặng ở các chương sau.
    \item Khi $\tau > 0$: Cửa sổ quan sát bị co ngắn lại vì phần đuôi $n+\tau \ge N$ vượt ra ngoài phạm vi dữ liệu của khung. Số lượng cặp mẫu chồng lấn giảm tuyến tính xuống còn đúng $N - \tau = 638 - \tau\text{ cặp}$.
\end{itemize}

\subsubsection*{Phép đếm số lượng phép tính MAC từ dãy số học}
Để xác định cao độ trong dải giọng nói của con người từ $F_{\text{floor}} = 75\text{ Hz}$ đến $F_{\text{ceiling}} = 600\text{ Hz}$, thuật toán Praat định nghĩa dải quét trễ (lag scan range) là $\tau \in [\tau_{\min}, \tau_{\max}] = [26, 214]$. Số lượng giá trị lag cần khảo sát là:
\begin{equation}
N_{\text{lags}} = \tau_{\max} - \tau_{\min} + 1 = 214 - 26 + 1 = 189\text{ độ trễ}
\end{equation}

Tại mỗi độ trễ $\tau \in [26, 214]$, phần cứng cần thực hiện đúng $638 - \tau$ phép nhân và tích lũy (Multiply-Accumulate -- MAC). Tổng số phép tính MAC cần thiết để quét qua toàn bộ 189 độ trễ là tổng của một cấp số cộng:
\begin{align}
\sum_{\tau=26}^{214} (638 - \tau) &= (638 - 26) + (638 - 27) + \dots + (638 - 214) \nonumber \\
&= 612 + 611 + \dots + 424
\end{align}
Áp dụng công thức tính tổng cấp số cộng với số số hạng $n = 189$, số hạng đầu $u_1 = 612$ và số hạng cuối $u_n = 424$:
\begin{equation}
S_{\text{scan}} = \frac{189 \times (612 + 424)}{2} = \frac{189 \times 1.036}{2} = 189 \times 518 = \mathbf{97.902\text{ phép MAC}}
\end{equation}

Khi cộng thêm giá trị $R[0]$ ($\tau = 0$ với $638 - 0 = 638\text{ phép MAC}$), tổng số phép nhân-cộng mà lõi phần cứng phải thực thi trên mỗi khung âm thanh là:
\begin{equation}
N_{\text{MAC, total}} = S_{\text{scan}} + (638 - 0) = 97.902 + 638 = \mathbf{98.540\text{ phép MAC}}
\end{equation}
(Con số $98.540$ này đã được kiểm chứng số học độc lập trong nhật ký \path{verify_v3/ch07_mac_math.log}).

\begin{thoughtqa}
\textbf{Câu hỏi đối chiếu:} Ở các tài liệu phiên bản v1 trước đây, người ta thường nhắc đến các con số $98.280$ và $98.920\text{ MAC}$. Các con số đó từ đâu ra?
\par\noindent\textbf{Trả lời:} Trong phiên bản cũ v1, chiều dài cửa sổ được lấy tròn ngây thơ là $N = 640\text{ mẫu}$ ($40\text{ ms} \times 16\text{ kHz}$). Khi đó:
\begin{equation}
\sum_{\tau=26}^{214} (640 - \tau) = \frac{189 \times (614 + 426)}{2} = \frac{189 \times 1.040}{2} = \mathbf{98.280\text{ phép MAC}}
\end{equation}
Cộng thêm $R[0]$ gồm $640\text{ phép MAC}$ sẽ cho ra đúng $98.280 + 640 = \mathbf{98.920\text{ phép MAC}}$. Sự chênh lệch giữa v1 và v2 là $98.920 - 98.540 = 380\text{ phép MAC}$, xuất phát từ quyết định sửa đúng độ dài cửa sổ $N = 638\text{ mẫu}$ theo chuẩn Praat 7.0.03 (được quy định tại \path{fon/Sound_to_Pitch.cpp}, dòng 334--339). Tài liệu này sử dụng triệt để con số chuẩn tắc $98.540\text{ phép MAC}$.
\end{thoughtqa}

\subsubsection*{Tương đương toán học giữa MAC trực tiếp và FFT 1024 điểm của Praat}
Trong phần mềm Praat của Paul Boersma (\path{fon/Sound_to_Pitch.cpp}), hàm tự tương quan không được tính bằng cách chạy hai vòng lặp lồng nhau trong miền thời gian, mà được tính thông qua cặp biến đổi Fourier nhanh (FFT) độ dài $1024\text{ điểm}$. Các dòng mã nguồn được kiểm chứng xác thực tại phiên bản Praat 7.0.03 gồm:
\begin{itemize}[leftmargin=1.5em, itemsep=2pt]
    \item \textbf{Dòng 344:} Xác định độ trễ quét tối thiểu:
    \begin{lstlisting}[language=C++, numbers=none]
minimumLag = max (2, Melder_ifloor (1.0 / my dx / pitchCeiling));
    \end{lstlisting}
    Tại $f_s = 16.000\text{ Hz}$ và $F_{\text{ceiling}} = 600\text{ Hz}$, $1 / (dx \cdot 600) = 16.000 / 600 \approx 26{,}67$, hàm sàn nguyên trả về $\mathbf{\tau_{\min} = 26\text{ mẫu}}$.
    \item \textbf{Dòng 345:} Xác định độ trễ quét tối đa:
    \begin{lstlisting}[language=C++, numbers=none]
maximumLag = min (Melder_ifloor (nsamp_window / periodsPerWindow) + 2, nsamp_window);
    \end{lstlisting}
    Tại $N = 638$ và $\text{periodsPerWindow} = 3$, $638 / 3 \approx 212{,}67$, hàm sàn cho $212$, cộng thêm $2$ biên an toàn cho ra đúng $\mathbf{\tau_{\max} = 214\text{ mẫu}}$.
    \item \textbf{Dòng 401--403:} Tính kích thước biến đổi FFT:
    \begin{lstlisting}[language=C++, numbers=none]
nsampFFT = 1;
while (nsampFFT < nsamp_window * (1.0 + interpolation_depth))
    nsampFFT *= 2;
    \end{lstlisting}
    Với chiều sâu nội suy $\text{interpolation\_depth} = 0{,}5$, ta có $638 \times 1{,}5 = 957\text{ mẫu}$. Lũy thừa nguyên bậc hai nhỏ nhất lớn hơn hoặc bằng $957$ là $\mathbf{N_{\text{FFT}} = 1024\text{ điểm}}$.
    \item \textbf{Dòng 79--80:} Đệm thêm các mẫu $0$ vào đuôi khung âm thanh:
    \begin{lstlisting}[language=C++, numbers=none]
for (integer j = nsamp_window + 1; j <= nsampFFT; j ++)
    frame [j] = 0.0;
    \end{lstlisting}
    Vùng từ mẫu số $639$ đến $1024$ ($386\text{ mẫu}$) được gán bằng $0{,}0$.
    \item \textbf{Dòng 144--156:} Thực hiện tự tương quan qua miền tần số với đoạn chú thích nổi tiếng của Boersma:
    \begin{lstlisting}[language=C++, numbers=none]
/* The FFT of the autocorrelation is the power spectrum. */
NUMfft_forward (frame.asSlice ());
// Tinh pho cong suat ac[i] = re^2 + im^2 tai moi tan so
NUMfft_backward (frame.asSlice ());
    \end{lstlisting}
    Phép biến đổi FFT thuận đưa tín hiệu sang miền tần số $X[k]$; bình phương độ lớn tạo ra mật độ phổ công suất $P[k] = |X[k]|^2$; và phép biến đổi FFT nghịch (IFFT) chuyển phổ công suất trở lại miền thời gian để thu được hàm tự tương quan $ac[\tau]$.
    \item \textbf{Dòng 158--164:} Chuẩn hóa hàm tự tương quan:
    \begin{lstlisting}[language=C++, numbers=none]
r [0] = 1.0;
for (integer i = 1; i <= brent_ixmax; i ++)
    r [i] = ac [i + 1] / (ac [1] * windowR [i + 1]);
    \end{lstlisting}
    Chia giá trị tự tương quan thô cho tích giữa năng lượng khung $ac[1]$ ($R[0]$) và tự tương quan của chính cửa sổ $windowR[i+1]$ ($R_w[\tau]$).
\end{itemize}

\begin{figure}[H]
\centering
\resizebox{0.92\textwidth}{!}{% chuong_v3/fig_ch07_direct_mac_vs_fft.tex
% So sánh đường ống tính tự tương quan: MAC trực tiếp trên FPGA vs FFT 1024 điểm trong Praat

\begin{tikzpicture}[
    >=latex,
    font=\small,
    box/.style={draw=black!80, rounded corners=3pt, line width=0.8pt, fill=white, align=center, inner sep=5pt},
    databox/.style={draw=blue!70!black, rounded corners=3pt, line width=0.9pt, fill=blue!5, align=center, inner sep=5pt},
    fftbox/.style={draw=orange!80!black, rounded corners=3pt, line width=0.9pt, fill=orange!5, align=center, inner sep=5pt},
    lbl/.style={font=\scriptsize, text=black!80, fill=white, inner sep=2pt},
    titlelbl/.style={font=\bfseries\small, text=black!90}
]

    % =========================================================================
    % Phần trên: Tuyến (a) MAC trực tiếp trên FPGA
    % =========================================================================
    \node[titlelbl, anchor=west] at (-0.5, 3.8) {\textbf{(a) Tuyến FPGA: Tự tương quan bằng MAC trực tiếp}};

    \node[databox, minimum width=2.2cm, minimum height=1.35cm] (ram) at (0.7, 2.5) {
        \textbf{Bộ đệm $x_w$}\\[2pt]
        \scriptsize 638 mẫu $\times$ 17-bit\\
        \scriptsize (2 khối M10K)
    };

    \node[databox, minimum width=2.2cm, minimum height=1.35cm] (addr) at (4.2, 2.5) {
        \textbf{Phát địa chỉ}\\[2pt]
        \scriptsize Quét $(n, n+\tau)$\\
        \scriptsize 98.540 cặp mẫu
    };

    \node[databox, minimum width=2.6cm, minimum height=1.35cm] (dsp) at (8.1, 2.5) {
        \textbf{Khối MAC DSP}\\[2pt]
        \scriptsize 17$\times$17-bit $\to$ 48-bit\\
        \scriptsize 1 chu kỳ/MAC (1 DSP)
    };

    \node[databox, minimum width=2.4cm, minimum height=1.35cm] (out_r) at (12.7, 2.5) {
        \textbf{Chuỗi $R[\tau]$}\\[2pt]
        \scriptsize 190 giá trị\\
        \scriptsize $98.543\text{ chu kỳ (RTF } 0{,}197\text{)}$
    };

    \draw[->, line width=1.0pt, blue!70!black] (ram.east) -- (addr.west) 
        node[midway, above=3pt, lbl] {$2\text{ mẫu/nhịp}$};
    \draw[->, line width=1.0pt, blue!70!black] (addr.east) -- (dsp.west) 
        node[midway, above=3pt, lbl] {$(a_r, b_r)$};
    \draw[->, line width=1.0pt, blue!70!black] (dsp.east) -- (out_r.west) 
        node[midway, above=3pt, lbl] {$R[0], R[\tau]$};

    % =========================================================================
    % Phần dưới: Tuyến (b) FFT 1024 điểm của Praat
    % =========================================================================
    \node[titlelbl, anchor=west] at (-0.5, 1.0) {\textbf{(b) Tuyến Praat: Tự tương quan qua cặp biến đổi FFT 1024 điểm}};

    \node[fftbox, minimum width=1.9cm, minimum height=1.3cm] (pad) at (0.7, -0.4) {
        \textbf{Đệm số 0}\\[2pt]
        \scriptsize $638 \to 1024\text{ mẫu}$\\
        \scriptsize Chống cuộn vòng
    };

    \node[fftbox, minimum width=2.1cm, minimum height=1.3cm] (fft1) at (3.5, -0.4) {
        \textbf{FFT 1024 điểm}\\[2pt]
        \scriptsize Bướm Radix-2/4\\
        \scriptsize Bảng twiddle ROM
    };

    \node[fftbox, minimum width=2.2cm, minimum height=1.3cm] (psd) at (6.7, -0.4) {
        \textbf{Phổ công suất}\\[2pt]
        \scriptsize $P[k] = |X[k]|^2$\\
        \scriptsize Quản lý tăng bit
    };

    \node[fftbox, minimum width=2.2cm, minimum height=1.3cm] (ifft) at (9.9, -0.4) {
        \textbf{IFFT 1024 điểm}\\[2pt]
        \scriptsize Biến đổi ngược\\
        \scriptsize Ra $ac[\tau]$ vòng
    };

    \node[fftbox, minimum width=2.2cm, minimum height=1.3cm] (cut) at (12.8, -0.4) {
        \textbf{Trích dải}\\[2pt]
        \scriptsize $\tau \le 214$\\
        \scriptsize Chia $R_w[\tau]$
    };

    \draw[->, line width=0.9pt, orange!80!black] (pad.east) -- (fft1.west);
    \draw[->, line width=0.9pt, orange!80!black] (fft1.east) -- (psd.west);
    \draw[->, line width=0.9pt, orange!80!black] (psd.east) -- (ifft.west);
    \draw[->, line width=0.9pt, orange!80!black] (ifft.east) -- (cut.west);

    % Nhãn so sánh dưới chân
    \node[draw=black!30, dashed, fill=black!2, rounded corners=3pt, inner sep=5pt, anchor=north] 
        at (6.8, -1.4) {
        \small\textbf{Kết luận kiến trúc:} Tại 50 MHz, MAC chỉ tiêu tốn 1 DSP và chiếm 19,7\% ngân sách ($98.543/500.000\text{ nhịp}$), tránh được độ phức tạp phần cứng của FFT.
    };

\end{tikzpicture}
}
\caption{So sánh quy trình tính toán tự tương quan: Tuyến FPGA dùng MAC trực tiếp và Tuyến Praat dùng cặp biến đổi FFT 1024 điểm.}
\label{fig:ch07_direct_mac_vs_fft}
\end{figure}

\textbf{Đọc hình.} Hình~\ref{fig:ch07_direct_mac_vs_fft} đặt cạnh nhau hai triết lý hiện thực hóa. Tuyến (a) trên FPGA triển khai luồng xử lý thẳng: các cặp mẫu được đọc trực tiếp từ bộ nhớ đệm M10K đưa vào bộ nhân-cộng DSP duy nhất, kết quả $R[\tau]$ được tích lũy tuần tự qua $98.543\text{ chu kỳ}$, hoàn toàn không cần bộ nhớ trung gian hay quản lý phức tạp. Ngược lại, Tuyến (b) của Praat trên máy tính sử dụng FFT 1024 điểm: tín hiệu buộc phải đệm số 0 lên 1024 mẫu, chuyển qua tầng bướm FFT phức, tính phổ công suất rồi IFFT ngược trở lại. Mặc dù hai cách đi khác nhau, chúng quy tụ về cùng một đáp số chính xác.

\begin{equationcard}
\textbf{Định lý triệt tiêu cuộn vòng thời gian (Zero Time-Domain Wrap-Around).}
\par\noindent Trong xử lý tín hiệu số, phép nhân trong miền tần số $P[k] = X[k] \cdot X^*[k]$ tương đương với \textbf{tự tương quan vòng} (Circular Autocorrelation) trong miền thời gian có chu kỳ lặp lại đúng bằng kích thước biến đổi $N_{\text{FFT}} = 1024$:
\begin{equation}
ac_{\text{circ}}[\tau] = \sum_{n=0}^{N_{\text{FFT}}-1} z[n] \cdot z[(n+\tau) \bmod N_{\text{FFT}}]
\end{equation}
trong đó $z[n]$ là chuỗi $x_w[n]$ được đệm số $0$: $z[n] = x_w[n]$ khi $0 \le n < N$, và $z[n] = 0$ khi $N \le n < N_{\text{FFT}}$.
\par\noindent Biểu thức $(n+\tau) \bmod N_{\text{FFT}}$ chỉ bị cuộn vòng vượt quá $1024$ khi $n + \tau \ge N_{\text{FFT}}$. Do $z[n] = 0$ với mọi $n \ge N = 638$, tích $z[n] \cdot z[n+\tau]$ chỉ khác $0$ khi $n \le 637$. Giá trị lớn nhất mà $n+\tau$ có thể đạt tới là:
\begin{equation}
(n+\tau)_{\max} = (N - 1) + \tau = 637 + \tau
\end{equation}
Để không xảy ra hiện tượng cuộn vòng ($n+\tau < N_{\text{FFT}}$), điều kiện cần và đủ là:
\begin{equation}
637 + \tau < 1024 \quad \Longleftrightarrow \quad \tau \le 1024 - 638 = \mathbf{386\text{ mẫu}}
\end{equation}
Vì thuật toán Praat chỉ quét lag tối đa đến $\tau_{\max} = 214\text{ mẫu}$, và $214 \le 386$ (vẫn còn dư biên an toàn $386 - 214 = 172\text{ mẫu}$), \textbf{hoàn toàn không có bất kỳ mẫu nào bị cuộn vòng đè lên nhau!}
\par\noindent Do đó, tự tương quan vòng của Praat đồng nhất $100\%$ với tự tương quan tuyến tính của bộ MAC trực tiếp:
\begin{equation}
ac[\tau+1]_{\text{Praat FFT}} \equiv R[\tau]_{\text{FPGA Direct MAC}}, \quad \forall \tau \in [0, 214]
\end{equation}
\end{equationcard}

Kết quả kiểm nghiệm độc lập trong kịch bản \path{verify_v3/ch07_fft_vs_mac.py} trên 5 khung tín hiệu chuẩn (sóng sin $150\text{ Hz}$, khung giọng nói sạch, khung giọng nói nhiễu $\text{SNR} = 5\text{ dB}$, và hai khung stress cực hạn) cho thấy sai số tương đối cực đại giữa FFT 1024 điểm và MAC trực tiếp chỉ đạt mức:
\begin{equation}
\frac{\max |ac[\tau+1] - R[\tau]|}{|R[0]|} \le 3{,}23 \times 10^{-16}
\end{equation}
Con số này xấp xỉ mức sai số làm tròn máy hữu hạn của kiểu dấu phẩy động 64-bit (\texttt{double}), chứng minh tuyệt đối sự tương đương toán học giữa hai phương pháp.

\textit{Tóm lại: Bản chất của tự tương quan là đếm số cặp mẫu trùng lặp qua độ trễ $\tau$. Dải quét 189 lag cùng $R[0]$ đòi hỏi chính xác $98.540\text{ phép MAC}$; phép đệm 386 số 0 trong FFT 1024 điểm của Praat đảm bảo tự tương quan vòng không bị cuộn chồng và đồng nhất hoàn toàn với tích vô hướng tuyến tính của MAC trực tiếp.}


\subsection*{7.2. Vì sao chọn MAC trực tiếp thay vì FFT trên FPGA}
\addcontentsline{toc}{subsection}{7.2. Vì sao chọn MAC trực tiếp thay vì FFT trên FPGA}

\subsubsection*{Vấn đề thực tế: Ảo tưởng về độ phức tạp thuật toán $O(N \log N)$ trên phần cứng}
Một phản xạ rất tự nhiên của bất kỳ kỹ sư phần mềm nào khi nhìn vào bài toán tự tương quan là mở sách giáo khoa xử lý tín hiệu và thốt lên: "Tại sao lại dùng $O(N^2)$ phép nhân trong miền thời gian khi ta có thể dùng FFT với độ phức tạp chỉ $O(N \log N)$?". Trên máy tính đa năng (CPU), lập luận này hoàn toàn chính xác: Praat chọn FFT 1024 điểm vì CPU x86 có bộ nhớ đệm L1/L2 khổng lồ, các tập lệnh vector AVX-512 xử lý số thực dấu phẩy động cực mạnh, và việc giảm số phép nhân giúp giảm đáng kể thời gian chiếm dụng CPU.

Tuy nhiên, khi bước chân vào thế giới phần cứng chuyên dụng trên chip FPGA, logic thiết kế bị chi phối bởi những ràng buộc hoàn toàn khác: diện tích silicon, băng thông bộ nhớ nội, độ phức tạp của đường truyền dữ liệu (datapath), và việc quản lý tăng trưởng bit lượng tử. Hãy cùng mổ xẻ chi tiết từng khía cạnh để thấy vì sao cách tiếp cận FFT trên FPGA lại là một sự đánh đổi bất lợi, và vì sao kiến trúc MAC trực tiếp là quyết định tối ưu.

\subsubsection*{Đếm số lượng phép toán: FFT ít hơn trên giấy, nhưng nặng trên silicon}
Hãy so sánh số lượng phép nhân thực tế giữa hai giải pháp:
\begin{itemize}[leftmargin=1.5em, itemsep=2pt]
    \item \textbf{Với biến đổi FFT 1024 điểm:}
    Kích thước $N_{\text{FFT}} = 1024 = 2^{10}$ đòi hỏi $10\text{ tầng bướm}$ (stages). Mỗi tầng gồm $1024 / 2 = 512\text{ bướm Cooley-Tukey}$. Một phép nhân bướm phức chuẩn $(a + jb)(c + jd) = (ac - bd) + j(ad + bc)$ tiêu tốn 4 phép nhân thực và 2 phép cộng thực (hoặc 3 phép nhân và 5 phép cộng theo thuật toán Karatsuba). Bỏ qua các phép nhân suy biến với $\pm 1$ và $\pm j$, phép biến đổi FFT thuận 1024 điểm đòi hỏi xấp xỉ:
    \begin{equation}
    10 \text{ tầng} \times 512 \text{ bướm} \times 4 \text{ nhân thực} \approx 20.480\text{ phép nhân thực}
    \end{equation}
    Tiếp đó, phép tính phổ công suất $P[k] = \text{Re}^2[k] + \text{Im}^2[k]$ trên 513 điểm tần số độc lập tiêu tốn thêm $2 \times 513 = 1.026\text{ phép nhân thực}$.
    Cuối cùng, phép biến đổi IFFT nghịch 1024 điểm đòi hỏi thêm một lượng phép tính tương đương: $\approx 20.480\text{ phép nhân thực}$.
    Tổng cộng cho cả quy trình tự tương quan qua miền tần số là:
    \begin{equation}
    N_{\text{mult, FFT}} \approx 20.480 + 1.026 + 20.480 \approx \mathbf{41.986\text{ phép nhân thực}}
    \end{equation}
    (Con số này nằm trong khoảng xấp xỉ $4 \times 10.240$ phép nhân thực mà tài liệu phân tích).
    \item \textbf{Với tự tương quan MAC trực tiếp:}
    Tổng số phép nhân cần tính toán như đã dẫn xuất ở Mục 7.1 là:
    \begin{equation}
    N_{\text{mult, MAC}} = \mathbf{98.540\text{ phép nhân thực}}
    \end{equation}
\end{itemize}

Nhìn thuần túy trên giấy tờ, FFT thực hiện khoảng $42.000\text{ phép nhân}$, ít hơn một nửa so với $98.540\text{ phép nhân}$ của MAC trực tiếp. Nhưng cái giá phần cứng phải trả để thực hiện được $42.000\text{ phép nhân}$ của FFT là vô cùng nặng nề:

\begin{enumerate}[leftmargin=1.5em, itemsep=3pt]
    \item \textbf{Độ phức tạp của bộ nhớ trung gian (Scratchpad Memory):}
    Trong thuật toán FFT, ngõ ra của tầng bướm này là ngõ vào của tầng bướm kế tiếp. Để xử lý dòng, hệ thống buộc phải triển khai cấu trúc đệm bóng bàn (ping-pong buffer) hoặc bộ đệm trễ hồi tiếp (Delay Feedback). Một khối FFT 1024 điểm dấu phẩy tĩnh cần lưu trữ ít nhất $2 \times 1024\text{ từ phức}$ (mỗi từ gồm phần thực và phần ảo $18\text{-bit}$), tiêu tốn dung lượng bộ nhớ:
    \begin{equation}
    2 \times 1024 \times (18 + 18)\text{ bit} = 73.728\text{ bit nhớ}
    \end{equation}
    Con số này đòi hỏi tiêu tốn ít nhất \textbf{8 khối M10K} nhúng ($10.240\text{ bit/khối}$) chỉ để làm bộ nhớ trung gian lưu các giá trị chuyển tiếp giữa các tầng bướm. Trong khi đó, bộ MAC trực tiếp chỉ cần đúng \textbf{2 khối M10K} để chứa khung mẫu ngõ vào $x_w$, và hoàn toàn \textbf{không cần thêm một bit nhớ trung gian nào}!
    \item \textbf{Độ phức tạp của bộ quay pha Twiddle Factor và bộ đảo bit:}
    FFT đòi hỏi một bảng tra ROM chứa các hệ số quay lượng giác $\cos(2\pi k / 1024)$ và $-\sin(2\pi k / 1024)$, cùng với một khối logic phức tạp để sinh địa chỉ đảo bit (bit-reversal address generator) nhằm tái sắp xếp lại chuỗi dữ liệu ngõ ra. Logic điều khiển này tiêu tốn hàng trăm khối logic ALM quý giá của Cyclone V.
    \item \textbf{Vấn đề tăng trưởng độ rộng bit (Bit Growth) và sai số lượng tử:}
    Mỗi tầng bướm của FFT có thể làm biên độ tín hiệu tăng lên tối đa một hệ số $\sqrt{2}$ (hoặc tăng 1 bit nhị phân trong trường hợp xấu nhất). Trải qua 10 tầng bướm, độ động của tín hiệu có thể phình to thêm tới 10 bit. Trên phần cứng FPGA dấu phẩy tĩnh, người thiết kế buộc phải chọn một trong ba phương án đầy rủi ro:
    \begin{itemize}[leftmargin=1.2em, itemsep=2pt]
        \item \emph{Dịch tỉ lệ cố định (Fixed Scaling):} Dịch phải 1 bit sau mỗi tầng bướm để chống tràn số. Hậu quả là các thành phần tín hiệu nhỏ bị chặt cụt liên tục, làm tích lũy nhiễu lượng tử trầm trọng và làm biến dạng các đỉnh tự tương quan ở các khung âm lượng nhỏ.
        \item \emph{Dấu phẩy động khối (Block Floating Point -- BFP):} Dò tìm đỉnh biên độ cực đại sau mỗi tầng để quyết định dịch tỉ lệ có điều kiện, đồng thời duy trì một số mũ chung. Giải pháp này đòi hỏi mạch dò biên độ tổ hợp phức tạp và làm nghẽn đường ống xung nhịp.
        \item \emph{Mở rộng datapath lên 32-bit:} Làm tăng gấp đôi tài nguyên bộ nhân DSP (buộc phải ghép nhiều bộ nhân $18 \times 18$ hoặc dùng DSP $27 \times 27$).
    \end{itemize}
    Ngược lại hoàn toàn, bộ MAC trực tiếp hoạt động trên thanh ghi tích lũy $48\text{-bit}$ hoàn toàn không nén, không dịch bit trung gian: mỗi phép nhân $17 \times 17\text{-bit}$ sinh ra tích số $34\text{-bit}$ chính xác tuyệt đối và được cộng dồn trực tiếp vào thanh ghi $48\text{-bit}$. \textbf{Sai số lượng tử trung gian bằng đúng $0$!}
\end{enumerate}

\subsubsection*{Ngân sách thời gian thực: Khi "nhanh hơn" trở nên vô nghĩa}
Hãy đặt con số $98.540\text{ phép MAC}$ vào bài toán thời gian thực cụ thể trên kit DE10-Standard. Tần số xung nhịp hệ thống trên FPGA là $f_{\text{clk}} = 50\text{ MHz}$ (chu kỳ xung nhịp $T_{\text{clk}} = 20\text{ ns}$). Tần số lấy mẫu âm thanh là $f_s = 16.000\text{ Hz}$. 

Mỗi chu kỳ lấy mẫu âm thanh tương ứng với số chu kỳ xung nhịp FPGA là:
\begin{equation}
N_{\text{clk/sample}} = \frac{50.000.000\text{ Hz}}{16.000\text{ Hz}} = \mathbf{3.125\text{ chu kỳ}}
\end{equation}

Thuật toán Praat quy định khoảng cách bước nhảy giữa hai khung liên tiếp (hop size) là $10\text{ ms}$, tương ứng với đúng $H = 160\text{ mẫu}$ âm thanh (được xác định tại dòng 296: $\text{timeStep} = \text{periodsPerWindow} / \text{pitchFloor} / 4 = 3 / 75 / 4 = 0{,}01\text{ s}$). Ngân sách thời gian thực tối đa mà FPGA có được giữa hai khung âm thanh liên tiếp là:
\begin{equation}
N_{\text{budget/hop}} = H \times N_{\text{clk/sample}} = 160 \times 3.125 = \mathbf{500.000\text{ chu kỳ xung nhịp}}
\end{equation}

Một bộ MAC phần cứng đơn lẻ trên Cyclone V thực thi với thông lượng \textbf{1 phép MAC trên 1 chu kỳ xung nhịp}. Với tổng số $98.540\text{ cặp toán hạng}$ cần tính toán cộng thêm 3 chu kỳ trễ đường ống, toàn bộ quá trình tính toán tự tương quan trên một khung $638\text{ mẫu}$ hoàn thành trong đúng:
\begin{equation}
N_{\text{cycles/frame}} = 98.540 + 3 = \mathbf{98.543\text{ chu kỳ xung nhịp}}
\end{equation}

Thời gian thực tế để tính xong một khung âm thanh là:
\begin{equation}
T_{\text{calc}} = 98.543 \times 20\text{ ns} \approx 1{,}97086\text{ ms}
\end{equation}

Hệ số thời gian thực (Real-Time Factor -- RTF) của riêng khối MAC được xác định bằng:
\begin{equation}
\text{RTF}_{\text{MAC}} = \frac{N_{\text{cycles/frame}}}{N_{\text{budget/hop}}} = \frac{98.543}{500.000} \approx \mathbf{0{,}197086 \approx 19{,}7\%}
\end{equation}

\begin{physicalinsight}
\textbf{Nhìn thấy vật lý: Sự xa xỉ của ngân sách thời gian thực.}
Khối MAC trực tiếp chỉ tiêu tốn đúng \textbf{$19{,}7\%$ ngân sách thời gian thực} của một bước nhảy khung $10\text{ ms}$! Sau khi tính xong trong chưa đầy $2\text{ ms}$, phần cứng MAC hoàn toàn rảnh rỗi trong hơn $8\text{ ms}$ còn lại ($401.457\text{ chu kỳ}$ nhàn rỗi) trước khi khung âm thanh tiếp theo xuất hiện.
\par\noindent Việc cố gắng nhồi nhét một cấu trúc FFT phức tạp để giảm thời gian tính từ $1{,}97\text{ ms}$ xuống $0{,}8\text{ ms}$ là hoàn toàn vô nghĩa trong bối cảnh này: nó không đem lại bất kỳ cải thiện nào cho người dùng cuối (vẫn phải chờ đủ $10\text{ ms}$ để có bước nhảy âm thanh mới), trong khi làm tiêu tốn thêm ít nhất 6 khối M10K và hàng trăm ALM, đồng thời đưa vào hệ thống nguy cơ sai số lượng tử tích lũy.
\end{physicalinsight}

Vì tất cả các lý do kỹ thuật trên, thiết kế này dứt khoát lựa chọn \textbf{kiến trúc MAC trực tiếp trong miền thời gian}. 

\textit{Tóm lại: FFT có vẻ ít phép nhân hơn trên lý thuyết, nhưng đòi hỏi bộ nhớ đệm M10K gấp $4\times$, logic điều khiển bướm phức tạp và dễ phát sinh méo lượng tử; trong khi đó, 1 bộ MAC trực tiếp chỉ chiếm $19{,}7\%$ ngân sách thời gian thực (RTF $\approx 0{,}197$) ở 50 MHz, bảo đảm độ chính xác số học nguyên bản và tiết kiệm tối đa tài nguyên silicon.}
